\section{Finitary relational presentations}
\label{sec:relational}

\subsection{Why single-valued names are sometimes the wrong interface}

Many constructive failures arise because a theorem is asked to return one
exactly selected representative when the mathematics only determines a finite
family of approximations.

\begin{itemize}
\item A polynomial with colliding roots has an unambiguous root multiset but no
      stable ordering of individual roots.
\item A real rotation near the identity has stable quaternion or approximate
      axis data, but no continuous choice of a unit axis at the identity.
\item A minimization problem may have several nearly minimizing grid points,
      even though the minimum value is stable.
\item Local analytic charts can overlap, while an exact chart selector at a
      boundary is discontinuous.
\item A finite proof may establish a compatible witness at every requested
      depth without supplying one globally chosen infinite branch.
\end{itemize}

Forcing a function in these situations either adds a separation hypothesis,
imports a choice principle, or makes the theorem false as a uniform
computational statement.  The framework therefore treats relations as a
first-class output type.

\subsection{Definition of a finitary relational presentation}

\begin{definition}[Finitary relational presentation]
A finitary relational presentation \(X\) consists of fixed
primitive-recursive data:
\begin{enumerate}
\item a parameter-domain Boolean \(\Dom_X(p)\);
\item a challenge-domain Boolean \(\mathsf{Chal}_X(p,q)\);
\item a finite answer bound \(B_X(p,q)\);
\item an answer verifier \(\Good_X(p,q,a)\in\B\), with the convention that
      valid answers satisfy \(a<B_X(p,q)\);
\item a refinement or compatibility verifier
      \(\Compat_X(p,q,a,q',a')\in\B\); and
\item optional observation and error decoders used to interpret an answer.
\end{enumerate}
All data are fixed symbols, not arbitrary run-time programs.
\end{definition}

The challenge \(q\) need not be a scalar precision.  It may encode a region,
derivative order, finite test function, matrix entry, truncation length, or a
bundle of such requests.  This makes the interface broad enough for analysis,
probability, topology, and algebra without pretending that every object is
merely a Cauchy sequence.

A parameter \(p\) is \emph{finitely inhabited to depth \(h\)} when every
challenge in the finite schedule up to \(h\) has an acceptable answer and the
answers can be connected by the required compatibility edges.  The exact
schedule is theorem-specific and primitive recursive.

\subsection{Functional presentations as a special case}

A usual fixed evaluator \(a(p,q)\) gives a relational presentation by setting
\[
 \Good_X(p,q,z)=1 \quad\Longleftrightarrow\quad z=a(p,q),
\]
and taking compatibility to be the desired regularity relation.  Thus no
existing functional development is lost.  The relational form is strictly
more flexible: it can represent finite root clusters, all valid pivots, all
near-minimizers, all cells intersecting a curve, or all charts whose guards
hold.

\subsection{Relational transducers}

\begin{definition}[Relational transducer]
A finitary relational transducer \(F:X\rightsquigarrow Y\) consists of:
\begin{enumerate}
\item a primitive-recursive finite query strategy.  Given an input parameter
      \(p\), output challenge \(r\), and current finite state, it either asks
      for an \(X\)-answer to a challenge \(q\) or halts;
\item a primitive-recursive output-answer bound and candidate generator;
\item a Boolean trace invariant and acceptance predicate; and
\item an embedded soundness derivation saying that any finite family of valid,
      mutually compatible input answers accepted by the trace yields a valid
      \(Y\)-answer.
\end{enumerate}
\end{definition}

The transducer need not choose among all valid input answers.  Its theorem is
relational: every accepted compatible input package leads to an acceptable
output package, or it produces a finite set of output candidates with a
verified nonemptiness and compatibility relation.

\subsection{Composition without hidden choice}

Suppose \(F:X\rightsquigarrow Y\) and \(G:Y\rightsquigarrow Z\).  A naive
composition would select one \(Y\)-answer for each query made by \(G\).  That
selection can destroy compatibility.  The relational composition instead
records a finite family of \(Y\)-candidates produced by \(F\), together with
their compatibility graph, and asks whether a compatible subfamily satisfies
\(G\)'s finite demand.  This is a finite constraint problem and is solved by
the viability dynamic program of \cref{sec:amalgamation}.

\begin{theorem}[Relational composition]
\label{thm:rel-composition}
The composition of two finitary relational transducers is a finitary
relational transducer.  Its verifier and finite candidate bounds are primitive
recursive.  Its soundness derivation is obtained by the embedded trace
compiler and the finite witness-amalgamation theorem.
\end{theorem}

\begin{proof}
Flatten the two traces into one phase-indexed trace as in
\cref{thm:trace-composition}.  Whenever \(G\) requests a \(Y\)-answer, invoke
\(F\)'s finite candidate generator rather than an exact selector.  The
composite state stores the finite candidate levels and compatibility edges.
After all requests are known, run backward viability to mark candidates that
can participate in a complete compatible assignment.  If at least one root
candidate is viable, bounded search returns a compatible assignment; if the
output is itself intended to remain relational, return all viable output
candidates.  Soundness follows by first applying the viability theorem to get
the finite compatible family, then applying the soundness derivations of
\(F\) and \(G\) by substitution and cut.  Every loop is bounded by the supplied
finite candidate bounds.
\end{proof}

\subsection{Three levels of output strength}

A theorem should state the weakest output that supports its applications.

\begin{description}[leftmargin=3.5cm,style=nextline]
\item[Relational output.] For every requested finite challenge, return a
      nonempty finite set of correct answers plus compatibility information.
      This needs no global selector.
\item[Approximate functional output.] Return one answer at each precision,
      with a regularity theorem.  This requires a shrinking or determinacy
      modulus as in \cref{sec:shrinking}.
\item[Exact selected output.] Return a distinguished exact representative,
      such as an ordered root, exact sign, normalized eigenvector, or branch.
      This requires a canonical algebraic normalization or a positive
      separation promise.  Without it the request may be obstructed.
\end{description}

The relational form is not a consolation prize.  Many classical theorems are
naturally invariant under permutations, signs, projective scaling, or chart
choices.  Stating the theorem at that invariant level often removes a
constructive obstruction while preserving all mathematically meaningful
content.

\subsection{Challenge trees rather than fixed tables}

An adaptive algorithm does not know its full input demand in advance.  Its
finite demand is a tree: a rational comparison at one stage selects the next
query.  A certificate can include one realized branch, while a uniform
soundness theorem verifies every branch.  The trace state contains the branch
history, and the invariant records that every answer on the history is valid.

If the theorem must be robust under approximate comparisons, branch guards
must overlap.  For example, given an approximation \(r\) to \(x\) within
\(\varepsilon\), one branch may be certified for
\(x\le c+2\varepsilon\) and the other for \(x\ge c-2\varepsilon\).  The
algorithm branches on whether \(r\le c\), but both mathematical branches are
valid in the overlap.  Exact branching is allowed only when an apartness gap
separates \(x\) from \(c\).

\subsection{Extensional interpretation}

A relational presentation represents mathematical content through the family
of accepted finite observations, not through the raw parameter code.  Two
parameters are observationally equivalent when every accepted answer for one
can be matched, at the requested precision, by an accepted answer for the
other and conversely.  This is a finite Hausdorff-style relation at each
precision.  Equality of represented objects is the universal closure over
precisions.

Operations should therefore be proved \emph{extensional}: matched input
observations at the precision requested by a modulus produce matched output
observations.  \Cref{sec:extensionality} gives the generic proof pattern.
